\scanpage{170}
% archive foliation: 169
de groupe un groupe constant $\uline{\pi_S}$, alors \unsure{tte} \uline{composante connexe} $P^\circ$ se surjecte sur $S$. Or c'est bien, car $P$ est \uline{trivial} !

Ceci étant, comme $P^{(0)}\subset\bigcap P^{(i)}$ pour \unsure{tt} $i$, on voit que $\bigcap_i P^{(i)}$ contient un pt de chaque fibre. \deleted{\unsure{En effet l'intersection}} \deleted{Il s'ensuit que} \deleted{\unsure{puisque}} \deleted{\unsure{les}} \inserted{chaque} $P^{(i)}$ est une réunion de composantes connexes de $P$, il en est de $=$ de l'intersection, \deleted{qui} $P'$, qui contient $\xi^\circ$. Elle est stable par $\pi'$, et on constate aussitôt (en regardant sur chaque fibre géométrique, et utilisant que $P'_{\bar s}\ne\varnothing$) que $P'$ est fibré principal homogène sous $\pi'$.
cqfd.

\uline{Conséquences} Si $\pi$ varie dans \inserted{une sous-}catégorie des groupes artiniens, \deleted{\unsure{et si}} stable par produits finis et sous-groupes, ou si $\pi$ varie ds n'importe quelle catégorie stable par ces opérations, mais $S$ étant loc. noeth --- alors le foncteur $\pi\mapsto\Pi^1(S,\xi;\pi)$ est strictement pro-représentable.

\margin{\unsure{On aura} pro-groupe $\widetilde\pi_1(S,\xi)$ [qui ne peut plus être interprété comme un simple groupe topologique].}
